\section{Задача и доступные вычисления}
В минимуме на множестве норма градиента может не обращаться в ноль.
Полный зазор метода Франк-Вульф учитывает допустимые направления
и в минимуме равен нулю. Поэтому мы включаем зазор в условие
кривизны, выводим допустимый шаг спуска, а затем используем
квазар-выпуклость для оценки ошибки по функции.

Пусть $Q\subset\mathbb R^n$ является непустым компактным выпуклым
множеством, $\operatorname{diam}Q\le D$, $D>0$. Функция $f$
дважды непрерывно дифференцируема в окрестности $Q$.
Зафиксируем её минимизатор $x^*\in Q$. Нормы векторов евклидовы,
нормы матриц спектральны. Доступны значения функции, её градиенты
и точная минимизация линейной функции по всему $Q$.
В текущей точке $x$ вычисляем
\begin{equation}\label{eq:gap}
 s\in\arg\min_{v\in Q}\langle\nabla f(x),v\rangle,\qquad
 d=s-x,\qquad
 G(x)=\langle\nabla f(x),x-s\rangle
     =\max_{v\in Q}\langle\nabla f(x),x-v\rangle.
\end{equation}
Предположим, что
\begin{equation}\label{eq:gap_smoothness}
 \|\nabla^2f(x)\|\le L_0+L_1G(x)\quad(x\in Q),
 \qquad L_0,L_1\ge0.
\end{equation}
Это условие зависит от допустимого множества. В частности,
$L_1D^2$ является безразмерной величиной; множитель $D$ в оценке
роста зазора ниже необходим.

Для сходимости к глобальному минимуму используем
$\gamma$-квазар-выпуклость относительно $x^*$:
\begin{equation}\label{eq:quasar}
 \gamma\bigl(f(x)-f(x^*)\bigr)
       \le\langle\nabla f(x),x-x^*\rangle\quad(x\in Q),
 \qquad 0<\gamma\le1.
\end{equation}
Для выпуклой функции подходит $\gamma=1$. Это условие допускает
и невыпуклые функции; пример приведён в разделе~\ref{sec:examples}.
Для шага нужны верные $L_0,L_1,D$, а для проверки точности по
функции дополнительно используется $\gamma$. Точка $x^*$ и
значение $f(x^*)$ алгоритму не передаются.

Необходимое условие минимума на множестве даёт $G(x^*)=0$,
в том числе при $\nabla f(x^*)\ne0$.
Из \eqref{eq:gap} и \eqref{eq:quasar} следует
\begin{equation}\label{eq:certificate}
 0\le f(x)-f(x^*)\le \frac{G(x)}{\gamma}.
\end{equation}
Зазор ограничивает ошибку по функции сверху. Обратная оценка
из одной лишь квазар-выпуклости не следует.

\subsection{Связь с литературой и границы вклада}
Полный зазор и оценка ошибки через него стандартны~\cite{jaggi}.
Модель кривизны через норму градиента и анализ Выгузова и
Стонякина~\cite{vyguzov}, опубликованные также совместно с
Гасниковым~\cite{vyguzovjournal}, мотивируют рассмотрение
\eqref{eq:gap_smoothness}. Мартинес-Рубио~\cite{quasar} уже изучает
ограниченную квазар-выпуклую оптимизацию методом Франк-Вульф;
теорема~9 даёт $6LD^2/((k+1)\gamma^2)$ при обычной гладкости.
Таким образом, квазар-выпуклый вариант FW и множитель $\gamma^{-2}$
не являются новыми результатами настоящей работы.

Используемый ниже накопленный сертификат тоже имеет известный
выпуклый аналог: при $\gamma=1$ он точно совпадает с best primal-dual gap
из раздела~2.2 работы Wirth, Pe\~na и Pokutta~\cite{wirth}.
Условия роста в этой работе связывают остаток линейной модели со
степенями $G$. В разделе~\ref{sec:structure} установлена явная связь
с нашим предположением. Мы применяем сертификат к анализу выбранного
шага при $0<\gamma\le1$, не заявляя новый выпуклый сертификат.

Для сравнения шагов существенны адаптивные альтернативы.
Backtracking FW оценивает кривизну проверками значений функции~\cite{pedregosa},
а auto-conditioned FW использует кривизну в пробной точке без
внутреннего поиска~\cite{acfw}. Расширения на conditional gradient
без поиска шага рассматриваются и в~\cite{yagishita}.
Адаптивные шаги Брегмана используют другую модель --- относительную
гладкость~\cite{bregman}. Эти предположения не отождествляются с
\eqref{eq:gap_smoothness}.

\begin{table}[ht]
\centering\small
\caption{Место настоящего анализа. Константы разных моделей имеют
разный смысл; таблица не ранжирует методы.}
\begin{tabular}{p{.23\linewidth}p{.34\linewidth}p{.34\linewidth}}
\toprule
Анализ & Данные о кривизне или росте & Существенное отличие\\
\midrule
Norm-based FW~\cite{vyguzov} & $K_0+K_1\|\nabla f\|$ & Явный короткий шаг через норму градиента\\
Квазар-выпуклый FW~\cite{quasar} & Общая константа $L$ & Известная оценка при ограничениях\\
Open-loop FW~\cite{wirth} & Остаток модели через $G^r$ & Известный накопленный выпуклый сертификат\\
Адаптивный FW~\cite{pedregosa,acfw} & Информация в пробных точках & Не требует заданных $L_0,L_1$\\
Настоящий анализ & $L_0+L_1G$ & Перенос кривизны, минимум модели и сертифицированная остановка\\
\bottomrule
\end{tabular}
\end{table}

Вклад работы состоит в выводе и анализе этой модели: лемме переноса
вдоль направления, логарифмическом минимизаторе модели, раздельных
вкладах $L_0/\varepsilon$ и $L_1\log(1/\varepsilon)$ и проверяемых
примерах. Показано также, что обычное правило Армихо сохраняет тот же
порядок числа обновлений. Явный шаг экономит пробные вычисления
функции при известных допустимых параметрах; всеобщее численное
превосходство над адаптивным FW не утверждается.
